\uuid{n8Xc}
\exo7id{3479}
\titre{exo7 3479}
\auteur{quercia}
\organisation{exo7}
\datecreate{2010-03-10}
\isIndication{false}
\isCorrection{true}
\chapitre{Déterminant, système linéaire}
\sousChapitre{Système linéaire, rang}
\cpge{mpsi-mp2i.calcul-matriciel-systemes.e, pcsi.determinants.c, ptsi.determinants.a, mpsi-mp2i.matrices-applications-lineaires.b, pcsi.matrices-applications-lineaires.b, ptsi.matrices-applications-lineaires.b, mpsi-mp2i.determinants.d, pcsi.calcul-matriciel-systemes-lineaires.d, ptsi.calcul-matriciel-systemes-lineaires.d, bcpst1.matrices, bcpst1.applications-lineaires-matrices}
\module{Algèbre}
\niveau{L2}
\difficulte{}

\contenu{
\question{
Soit $M \in \mathcal{M}_3(K)$ antisymétrique. Quel est le rang de $M$ ?
}
\reponse{
2 ou 0.
}
}
